\section{Strategic Projection and Equilibrium Transfer}

Let
\[
\Fix(G)=\{\widehat\Gamma:g\widehat\Gamma=\widehat\Gamma\ \text{for all }g\in G\}
\]
be the linear subspace of exactly invariant games.

\begin{definition}[Strategic symmetry defect]
The defect of a game relative to a candidate group is
\begin{equation}
\delta_G(\Gamma)
=
\min_{\widehat\Gamma\in\Fix(G)}
\devdist(\Gamma,\widehat\Gamma).
\label{eq:defect}
\end{equation}
A minimizer is a nearest symmetric surrogate.
\end{definition}

A zero defect is weaker than literal payoff symmetry: it means that $\Gamma$ is strategically equivalent to an exactly symmetric game.  This is intentional because nonstrategic payoff components do not affect equilibrium behavior.

\subsection{Projection as a Linear Program}

For each surrogate payoff coordinate introduce $\widehat u_i(a)$; for each slice $(i,a_{-i})$ introduce lower and upper residual variables $L_{i,a_{-i}}$ and $U_{i,a_{-i}}$; and introduce $t\ge0$.  Consider
\begin{align}
\min\quad &t \label{eq:projection-lp}\\
\text{s.t.}\quad
&L_{i,a_{-i}}
\le u_i(a)-\widehat u_i(a)
\le U_{i,a_{-i}} &&\forall i,a,\nonumber\\
&U_{i,a_{-i}}-L_{i,a_{-i}}\le t
&&\forall i,a_{-i},\nonumber\\
&\widehat u_i(a)=\widehat u_{\pi_g(i)}(ga)
&&\forall g\in S,i,a.\nonumber
\end{align}
Only a generating set $S$ is needed because invariance under generators implies invariance under the generated group.

\begin{theorem}[Exact strategic projection]
\label{thm:lp}
The optimum of \eqref{eq:projection-lp} equals $\delta_G(\Gamma)$, and its payoff variables define a nearest exactly $G$-invariant surrogate.
\end{theorem}
\begin{proof}
For fixed $\widehat u$, the smallest feasible interval for a residual slice has width equal to the range in \eqref{eq:range-form}.  The largest such width is exactly $\devdist(\Gamma,\widehat\Gamma)$.  The equality constraints characterize $\Fix(G)$.
\end{proof}

If $P=n\prod_i|A_i|$ is the number of payoff entries and
$Q=\sum_i\prod_{j\ne i}|A_j|$ is the number of slices, the LP has
$P+2Q+1$ variables and $2P+Q$ strategic inequalities, plus sparse symmetry equalities.  Additional linear constraints can preserve a game class $\mathcal C$, such as zero-sum or common-payoff games.  We write
\begin{equation}
\delta_{G,\mathcal C}(\Gamma)
=
\min_{\widehat\Gamma\in\Fix(G)\cap\mathcal C}
\devdist(\Gamma,\widehat\Gamma)
\label{eq:structured-defect}
\end{equation}
for this structured projection.  The unrestricted and structured defects need not coincide.

\subsection{A Maximum-Mean-Cycle Characterization}

The unrestricted projection has a more explicit combinatorial form.  Let $V$ be the set of orbits of payoff coordinates $(i,a)$ under $G$, and denote the orbit of a coordinate by $[i,a]$.  Construct a directed multigraph $\mathcal D_G(\Gamma)=(V,E)$.  For every ordered unilateral comparison
$e=(i,a_{-i},b_i,a_i)$, add an edge
\[
[i,(b_i,a_{-i})]\longrightarrow[i,(a_i,a_{-i})]
\]
with label
\[
c_e=u_i(a_i,a_{-i})-u_i(b_i,a_{-i}).
\]
Both directions are present because all ordered pairs of distinct actions are included.  A $G$-invariant surrogate is specified by one potential $x_v$ for each orbit, with $\widehat u_i(a)=x_{[i,a]}$.

\begin{theorem}[Cycle characterization]
\label{thm:cycle}
Let $\mathcal Z$ be the set of directed cycles of $\mathcal D_G(\Gamma)$, including self-loops.  Then
\begin{equation}
\boxed{
\delta_G(\Gamma)
=
\max_{C\in\mathcal Z}
\left(-\frac{1}{|C|}\sum_{e\in C}c_e\right).
}
\label{eq:cycle-characterization}
\end{equation}
The maximum is defined as zero when $\mathcal Z$ is empty.  At the optimum, shortest-path potentials for edge lengths $c_e+\delta_G$ define a nearest invariant surrogate.  Hence $\delta_G$ and a witness cycle can be computed in polynomial time by maximum-cycle-mean and difference-constraint algorithms.
\end{theorem}
\begin{proof}
A potential $x$ has strategic error at most $t$ exactly when
$|x_{\mathrm{head}(e)}-x_{\mathrm{tail}(e)}-c_e|\le t$ for every comparison.  Because the reverse comparison is also an edge, this is equivalent to the one-sided difference constraints
\[
x_{\mathrm{head}(e)}-x_{\mathrm{tail}(e)}\le c_e+t
\quad\forall e.
\]
A system of difference constraints is feasible if and only if every directed cycle has nonnegative total length.  Thus feasibility is equivalent to
$\sum_{e\in C}(c_e+t)\ge0$ for every cycle $C$, or
$t\ge-|C|^{-1}\sum_{e\in C}c_e$.  Minimizing $t\ge0$ gives \eqref{eq:cycle-characterization}, including the acyclic case.  At the critical value there are no negative cycles, so shortest-path distances from an added source provide feasible potentials.  Karp's algorithm finds the extremal cycle mean in polynomial time \cite{karp1978cycle}.
\end{proof}

The witness cycle has a useful interpretation: exact symmetry identifies several payoff coordinates, while unilateral incentives prescribe differences between them.  A nonzero sum around a cycle is an inconsistent set of prescriptions.  The most-negative average inconsistency is precisely the minimum uniform relaxation required to make all prescriptions integrable.

\subsection{Equilibrium Preservation}

\begin{lemma}[Mixed incentive transfer]
\label{lem:mixed-transfer}
If $\devdist(\Gamma,\widehat\Gamma)\le\delta$, then for every player $i$, opponent mixture $\sigma_{-i}$, and own mixed strategies $p_i,q_i$,
\begin{align*}
\Big|&[u_i(p_i,\sigma_{-i})-u_i(q_i,\sigma_{-i})]\\
&-[\widehat u_i(p_i,\sigma_{-i})-\widehat u_i(q_i,\sigma_{-i})]\Big|
\le\delta.
\end{align*}
\end{lemma}
\begin{proof}
For each pure opponent profile, all residual values over player $i$'s actions lie in an interval of width at most $\delta$.  The difference between any two convex combinations of values in that interval is at most $\delta$; averaging over opponent profiles preserves the bound.
\end{proof}

\begin{theorem}[Equilibrium transfer]
\label{thm:transfer}
If $\sigma$ is an $\epsilon$-Nash equilibrium of $\widehat\Gamma$ and
$\devdist(\Gamma,\widehat\Gamma)\le\delta$, then $\sigma$ is an
$(\epsilon+\delta)$-Nash equilibrium of $\Gamma$.
\end{theorem}
\begin{proof}
For any unilateral deviation $\tau_i$, Lemma~\ref{lem:mixed-transfer} gives
\[
u_i(\tau_i,\sigma_{-i})-u_i(\sigma)
\le
\widehat u_i(\tau_i,\sigma_{-i})-\widehat u_i(\sigma)+\delta
\le\epsilon+\delta.
\]
\end{proof}

Every finite exactly $G$-invariant game has a Nash equilibrium respecting $G$ \cite{tewolde2025game}.  A self-contained proof, included in the supplement, restricts an equivariant logit-response map to the compact convex set of invariant profiles, applies Brouwer's theorem, and lets the response precision diverge.

\begin{corollary}[Invariant approximate equilibrium]
\label{cor:invariant-existence}
Every finite game $\Gamma$ admits a $G$-respecting
$\delta_G(\Gamma)$-Nash equilibrium.
\end{corollary}

The constant in this guarantee cannot be improved uniformly.
\begin{proposition}[Asymptotic tightness]
\label{prop:tight}
For every $k\ge2$, there is a two-player game and a cyclic action group $G_k$ with $\delta_{G_k}=1$ such that every $G_k$-respecting profile has one player's regret at least $1-1/k$.
\end{proposition}
\begin{proof}
Player 1 has $k$ actions, receives payoff one from one distinguished action and zero from the others, independently of player 2; player 2 has two irrelevant actions and zero payoff.  The group cycles player 1's actions.  Any invariant surrogate has equal payoffs across the orbit, so the defect is one.  The only invariant strategy is uniform, whose regret is $1-1/k$.
\end{proof}
